% GENERATED FILE. Edit canonical lessons, then run npm run generate:books.
% canonical-source-hash: fnv1a64:356c7d8ce52d6fe9
% canonical-lessons: GE-C305-feier, GE-C305-feierlich, GE-C305-feiertag, GE-C305-zeremonie, GE-C305-einweihen

\chapter{Die Feier}
\label{ch:ge-die-feier}
\hlchaptermodality{305}

\begin{chapteropening}
\textbf{By the end of this chapter:} \emph{I can talk about celebrations: a celebration, solemn, a public holiday, a ceremony and to open officially.}

Complete the last lesson of chapter 305: I can talk about celebrations: a celebration, solemn, a public holiday, a ceremony and to open officially.
\end{chapteropening}

\section[die Feier]{die Feier — the celebration, the party}

\label{lesson:GE-C305-feier}



\begin{quote}
Before the new one: say the German for the widow, then the German for to inherit.
\end{quote}

\subsection*{You'll want to know: die Feier}
\textbf{die Feier} — \textquotedblleft{}the celebration, the party\textquotedblright{}.

\begin{quote}
Die Feier ging bis Mitternacht.
\end{quote}

\begin{quote}
\emph{The celebration went on until midnight.}
\end{quote}

Plural: \emph{die Feiern}. The verb is \emph{feiern}.

\begin{cousinweb}
\emph{Feier} comes from Middle High German \emph{vire}, Old High German \emph{fīra}, which came from Latin \emph{fēriae} \textquotedblleft{}holidays, feast days\textquotedblright{} (EtymDB).
\end{cousinweb}

\begin{grammarlens}[title={What you've built}]
One more word for telling, explaining and describing at B1.
\end{grammarlens}

\subsection*{Your turn}
\begin{itemize}
\raggedright
  \item \emph{Say it:} \emph{die Feier}
  \item \emph{Say it:} \emph{die Feier}, once more
  \item \emph{Say it:} a sentence of your own with \emph{die Feier}
  \item \emph{Recall:} say the German for the widow, then the German for to inherit, then say \emph{die Feier} again
\end{itemize}

\begin{tcolorbox}[breakable,colback=teal!4,colframe=teal!35!black,title={Before you move on}]
What does die Feier mean? (\textquotedblleft{}The celebration, the party\textquotedblright{}.) Say it once more.
\end{tcolorbox}
\section[feierlich]{feierlich — solemn, festive}

\label{lesson:GE-C305-feierlich}



\begin{quote}
Before the new one: say the German for to inherit, then the German for the celebration.
\end{quote}

\subsection*{You'll want to know: feierlich}
\textbf{feierlich} — \textquotedblleft{}solemn, festive\textquotedblright{}.

\begin{quote}
Die Stimmung war sehr feierlich.
\end{quote}

\begin{quote}
\emph{The mood was very solemn.}
\end{quote}

\emph{Feierlich} is for a grand, serious occasion.

\begin{cousinweb}
\emph{feierlich} is built on \emph{Feier} \textquotedblleft{}celebration\textquotedblright{}, which came from Latin \emph{fēriae} \textquotedblleft{}holidays\textquotedblright{} (EtymDB), with the ending \emph{-lich}.
\end{cousinweb}

\begin{grammarlens}[title={What you've built}]
One more word for telling, explaining and describing at B1.
\end{grammarlens}

\subsection*{Your turn}
\begin{itemize}
\raggedright
  \item \emph{Say it:} \emph{feierlich}
  \item \emph{Say it:} \emph{feierlich}, once more
  \item \emph{Say it:} a sentence of your own with \emph{feierlich}
  \item \emph{Recall:} say the German for to inherit, then the German for the celebration, then say \emph{feierlich} again
\end{itemize}

\begin{tcolorbox}[breakable,colback=teal!4,colframe=teal!35!black,title={Before you move on}]
What does feierlich mean? (\textquotedblleft{}Solemn, festive\textquotedblright{}.) Say it once more.
\end{tcolorbox}
\section[der Feiertag]{der Feiertag — the public holiday}

\label{lesson:GE-C305-feiertag}



\begin{quote}
Before the new one: say the German for the celebration, then the German for solemn.
\end{quote}

\subsection*{You'll want to know: der Feiertag}
\textbf{der Feiertag} — \textquotedblleft{}the public holiday\textquotedblright{}.

\begin{quote}
Am Feiertag sind alle Geschäfte geschlossen.
\end{quote}

\begin{quote}
\emph{On the public holiday all the shops are closed.}
\end{quote}

Plural: \emph{die Feiertage}. \emph{Frohe Feiertage!} — happy holidays!

\begin{cousinweb}
\emph{Feiertag} joins \emph{Feier}, from Latin \emph{fēriae} \textquotedblleft{}holidays\textquotedblright{} (EtymDB), and the inherited \emph{Tag} \textquotedblleft{}day\textquotedblright{}.
\end{cousinweb}

\begin{grammarlens}[title={What you've built}]
One more word for telling, explaining and describing at B1.
\end{grammarlens}

\subsection*{Your turn}
\begin{itemize}
\raggedright
  \item \emph{Say it:} \emph{der Feiertag}
  \item \emph{Say it:} \emph{der Feiertag}, once more
  \item \emph{Say it:} a sentence of your own with \emph{der Feiertag}
  \item \emph{Recall:} say the German for the celebration, then the German for solemn, then say \emph{der Feiertag} again
\end{itemize}

\begin{tcolorbox}[breakable,colback=teal!4,colframe=teal!35!black,title={Before you move on}]
What does der Feiertag mean? (\textquotedblleft{}The public holiday\textquotedblright{}.) Say it once more.
\end{tcolorbox}
\section[die Zeremonie]{die Zeremonie — the ceremony}

\label{lesson:GE-C305-zeremonie}



\begin{quote}
Before the new one: say the German for solemn, then the German for the public holiday.
\end{quote}

\subsection*{You'll want to know: die Zeremonie}
\textbf{die Zeremonie} — \textquotedblleft{}the ceremony\textquotedblright{}.

\begin{quote}
Die Zeremonie dauerte eine Stunde.
\end{quote}

\begin{quote}
\emph{The ceremony lasted an hour.}
\end{quote}

Plural: \emph{die Zeremonien}. Stress at the end: \emph{Ze-re-mo-NIE}.

\begin{cousinweb}
\emph{Zeremonie} is a learned borrowing from Latin \emph{caerimōnia} \textquotedblleft{}a sacred rite\textquotedblright{} (Lewis \& Short); EtymDB lists no parent.
\end{cousinweb}

\begin{grammarlens}[title={What you've built}]
One more word for telling, explaining and describing at B1.
\end{grammarlens}

\subsection*{Your turn}
\begin{itemize}
\raggedright
  \item \emph{Say it:} \emph{die Zeremonie}
  \item \emph{Say it:} \emph{die Zeremonie}, once more
  \item \emph{Say it:} a sentence of your own with \emph{die Zeremonie}
  \item \emph{Recall:} say the German for solemn, then the German for the public holiday, then say \emph{die Zeremonie} again
\end{itemize}

\begin{tcolorbox}[breakable,colback=teal!4,colframe=teal!35!black,title={Before you move on}]
What does die Zeremonie mean? (\textquotedblleft{}The ceremony\textquotedblright{}.) Say it once more.
\end{tcolorbox}
\section[einweihen]{einweihen — to open officially, to inaugurate; to let in on}

\label{lesson:GE-C305-einweihen}



\begin{quote}
Before the new one: say the German for the public holiday, then the German for the ceremony.
\end{quote}

\subsection*{You'll want to know: einweihen}
\textbf{einweihen} — \textquotedblleft{}to open officially, to inaugurate; to let in on\textquotedblright{}.

\begin{quote}
Am Sonntag weihte die Stadt die neue Brücke ein.
\end{quote}

\begin{quote}
\emph{On Sunday the town officially opened the new bridge.}
\end{quote}

Separable: \emph{sie weihen ein}. Also to let someone in on a secret: \emph{jemanden einweihen}.

\begin{cousinweb}
\emph{einweihen} is built on \emph{weihen} \textquotedblleft{}to make holy\textquotedblright{}, Middle High German \emph{wīhen}, from Proto-Germanic \emph{wīhijaną} \textquotedblleft{}to hallow\textquotedblright{} (EtymDB).
\end{cousinweb}

\begin{grammarlens}[title={What you've built}]
One more word for telling, explaining and describing at B1.
\end{grammarlens}

\subsection*{Your turn}
\begin{itemize}
\raggedright
  \item \emph{Say it:} \emph{einweihen}
  \item \emph{Say it:} \emph{einweihen}, once more
  \item \emph{Say it:} a sentence of your own with \emph{einweihen}
  \item \emph{Recall:} say the German for the public holiday, then the German for the ceremony, then say \emph{einweihen} again
\end{itemize}

\begin{tcolorbox}[breakable,colback=teal!4,colframe=teal!35!black,title={Before you move on}]
What does einweihen mean? (\textquotedblleft{}To open officially, to inaugurate; to let in on\textquotedblright{}.) Say it once more.
\end{tcolorbox}
